\documentclass[10pt]{beamer}

%\usepackage[backend=bibtex,firstinits=true,style=verbose-inote,citestyle=authortitle]{biblatex}
\usepackage{bm}
\usepackage{graphicx}
\usepackage{subcaption}
\usepackage{amsmath}
\usepackage{makecell}
\usepackage{filecontents}
\usepackage{biblatex}
\newcommand{\expect}[2][]{
\ifthenelse{\equal{#1}{}}{
\mathbb{E}\left[#2\right]
}{
\underset{#1}{\mathbb{E}}\left[#2\right]
}}

\newcommand{\cov}[2][]{
\ifthenelse{\equal{#1}{}}{
\text{Cov}\left[#2\right]
}{
\underset{#1}{\text{Cov}}\left[#2\right]
}}


\newcommand{\var}[2][]{
\ifthenelse{\equal{#1}{}}{
\text{Var}[#2]
}{
\underset{#1}{\text{Var}}[#2]
}}

\newcommand{\loss}[2][]{
\ifthenelse{\equal{#1}{}}{
\mathcal{L}(#2)
}{
\mathcal{L}_{#1}(#2)
}}

\newcommand{\kl}[2]{
\text{D}_\text{KL}[#1 \parallel #2]
}

\newcommand{\R}{\mathbb{R}}
%\newcommand{\Prob}{\mathbb{P}}

\newcommand{\1}[1]{\mathds{1}\{#1\}}


%\usecolortheme{dolphin}
\setbeamertemplate{navigation symbols}{}
\setbeamertemplate{section in toc}{\inserttocsectionnumber.~\inserttocsection}


\title{Your Classifier is Secretly an Energy Based Model and You Should Treat it Like One}
%\subtitle{}
%\author{Ivan Skorokhodov}
%\date{}
%\logo{\includegraphics[height=1cm]{images/ipavlov-logo.png}}

\newcommand{\citepaper}[1]{\citetitle{#1} by \citeauthor{#1}}

%\graphicspath{{./images}}

%\usetheme{lucid}
\begin{document}

\begin{frame}
    \titlepage
\end{frame}

\begin{frame}
    \frametitle{What is an Energy Based Model (EBM)?}
\pause
Any pdf $p_\theta(x)$ can be written down as:
    
    \begin{equation}
p_{\theta}(\mathbf{x})=\frac{\exp \left(-E_{\theta}(\mathbf{x})\right)}{Z_\theta}
\end{equation}
\pause
where $Z_\theta$ is a normalizing constant:
\begin{equation*}
Z_\theta = \int_{\bm x} \exp{(-E_\theta(\bm x))}d\bm x
\end{equation*}

\begin{itemize}
    \item\pause This parametrization allows us to forget about $p_\theta(x)$ and work with $E_\theta(x)$
    \item\pause Function $E_\theta(x)$ is an \textit{energy function} (hence the name --- EBM)
    \item\pause Energy functions are \textit{easier} to work with: you do not have any restrictions on $E_\theta(x)$ ($p_\theta(x)$ is nonnegative and integrates to 1).
    \item\pause Energy functions are \textit{harder} to work with: $Z_\theta$ is hard or impossible to compute, so we can't optimize EBM for $\theta$
\end{itemize}
\end{frame}

\begin{frame}
    \frametitle{How to train EBMs}
    
    \begin{itemize}
        \item\pause We perform standard MLE: maximize $\log p_\theta(x_i)$ at our points $x_i$.
        \item\pause By optimizing $\log p_\theta(x)$ we also optimize $E_\theta(x)$.
        \item\pause We can't compute $\log p_\theta(x)$, but we can (approximately) compute its log-derivative via:
\begin{equation}
\frac{\partial \log p_{\theta}(\mathbf{x})}{\partial \theta}=\mathbb{E}_{p_{\theta}\left(\mathbf{x}^{\prime}\right)}\left[\frac{\partial E_{\theta}\left(\mathbf{x}^{\prime}\right)}{\partial \theta}\right]-\frac{\partial E_{\theta}(\mathbf{x})}{\partial \theta}
\end{equation}
        \item\pause Proof: follows directly from the definition.
        \item\pause But a new problem arises: how to compute expectation?
        \item\pause Answer: we can't do that directly, so let's use SGLD sampling.
    \end{itemize}
\end{frame}

\begin{frame}    
    \frametitle{Stochastic Gradient Langevin Dynamics (SGLD)}
    Motivation:
    \begin{itemize}
        \item\pause Imagine we have a ``magic'' function $p_\theta(x)$ that takes an input $x$ and tells us its density.
        \item\pause Unfortunately, we cannot sample new $x$-s from it directly, i.e. $p_\theta(x)$ does not support such operations.
        \item\pause But we really want to sample from it! What should we do?
        \item\pause SGLD is a method to sample from such functions.
    \end{itemize}
\end{frame}

\begin{frame}
    \frametitle{Stochastic Gradient Langevin Dynamics (SGLD)}
    SGLD in short:
    \begin{enumerate}
        \item\pause Start from any random point $x_0$.
        \item\pause Compute the gradient of $\nabla \log p_\theta(x_0)$.
        \item\pause Perform a small step in that direction.
        \item\pause Add a little bit of noise.
        \item\pause Obtain $x_{t+1}$.
    \end{enumerate}
    \pause In other words, run the following recurrence a lot of times:
\begin{equation}
\mathbf{x}_{0} \sim p_{0}(\mathbf{x}), \quad \mathbf{x}_{i+1}=\mathbf{x}_{i}-\frac{\alpha}{2} \frac{\partial E_{\theta}\left(\mathbf{x}_{i}\right)}{\partial \mathbf{x}_{i}}+\epsilon, \quad \epsilon \sim \mathcal{N}(0, \alpha)
\end{equation}

\pause Justification: if your model and your choice of $\alpha$ are ``good'' (in some sense), then after infinite amount of steps your samples will be ``good'' (in some sense).

\end{frame}

\begin{frame}
    \frametitle{Ok, how to treat a classifier as an EBM?}
    
    
    \begin{itemize}
        \item\pause Remember that \textit{any} function can be used as $E_\theta(x)$
        \item\pause Let $f_\theta(x)$ be our neural network.
        \item\pause Let $f_\theta(x)[y]$ be the $y$-th logit (i.e. logit value for class $y$).
        \item\pause Let our examples be $s = (x,y)$
        \item\pause We define energy $E_\theta(s) = E_\theta(x,y) = -f_\theta(x)[y]$.
        \item\pause Using $E_\theta(x,y)$ we can define $p_\theta(x)$:
        \begin{equation}
p_{\theta}(\mathbf{x})=\sum_{y} p_{\theta}(\mathbf{x}, y)=\frac{1}{Z_\theta}\sum_{y} \exp \left(f_{\theta}(\mathbf{x})[y]\right)
\end{equation}
        \item\pause Using $p_\theta(x)$ we can define $E_\theta(x)$:
        \begin{equation}
E_{\theta}(\mathbf{x})=-\log \sum_{y} \exp \left(f_{\theta}(\mathbf{x})[y]\right)
\end{equation}
        \item\pause Higher the logits --- lower the energy of $x$.
        \item\pause Our final loss looks like:
        \begin{equation}
            \mathcal{L}(x) = \log p(x,y) = \underbrace{\log p(y|x)}_\text{classification loss} + \underbrace{\log p(x)}_\text{generative loss}
        \end{equation}
    \end{itemize}
\end{frame}

\begin{frame}
    \frametitle{Results and discussion}
    \begin{itemize}
        \item\pause Disclaimer: results are really good.
        \item\pause Authors built a model on SVHN/CIFAR10/CIFAR100 which:
        \begin{itemize}
            \item\pause has good classification accuracy
            \item\pause has good IS/FID scores (i.e. has good samples)
            \item\pause has much better calibrated predictions (i.e. is not overconfident unlike traditional classifiers)
            \item\pause has good out-of-distribution predictions (i.e. detecting when we feed OOD-samples into it)
            \item\pause is robust to adversarial attacks
            \item\pause ...and all of these at the same time!
        \end{itemize}
        \item\pause A huge limitation is that because of SGLD the training is unstable and the model is hard to scale.
        \item\pause A big limitation: we cannot compute the loss value!!!
        \item\pause An approach was coined Joint Energy based Model (JEM) since we model $(x,y)$ simultaneously.
    \end{itemize}
\end{frame}

\end{document}
